\begin{abstract}
% Lift / refine from progress.md §1.
EEG foundation models (FMs) promise transferable representations for clinical
EEG, but their utility on weak-signal resting-state tasks remains unclear.
We conducted a systematic cross-subject evaluation of three architecturally
distinct FMs (LaBraM, REVE, CBraMod) on chronic stress classification (UCSD
Stress, $N{=}17$ subjects), and validated the diagnosis on two clinical
reference datasets (ADFTD Alzheimer's, $N{=}65$; TDBRAIN MDD, $N{=}359$).
Three findings emerge.
\textbf{(1)}~Frozen FM features encode subject identity 11--13$\times$ more
than diagnostic labels (variance decomposition via $\eta^{2}$), and this
pattern is consistent across architectures and tasks.
\textbf{(2)}~Trial-level cross-validation inflates reported balanced
accuracy by $\sim$21~points relative to subject-level cross-validation across
all three FMs, explaining the gap between published 90\% claims and our
55--66\% subject-level results.
\textbf{(3)}~A random forest on hand-crafted band-power features matches
fine-tuned LaBraM ($0.666$ vs.\ $0.656$~\BA{} on stress; $0.753$ vs.\ $0.752$
on Alzheimer's), and adversarial subject removal degrades rather than
improves performance.
We characterise the resting-state stress signal as method-independent at
$\sim$0.66~\BA{} under binary cross-subject framing, and argue that
within-subject longitudinal modelling --- where subject identity becomes the
reference frame rather than the noise --- is the productive path forward.
\end{abstract}
